\section*{Conclusões}

\rasc{No método dos painéis, $N = 260$ painéis bastam para o $c_l$ a $5^\circ$
ficar a menos de 0,5\,\% do valor obtido com 800 painéis. No CFD escolheu-se a
malha com cerca de 64 mil células, para a qual a incerteza relativa do $c_l$
entre malhas é de 1,3\,\%.

Os três métodos dão um ângulo de sustentação nula pequeno e negativo, entre
$-0{,}67^\circ$ e $-0{,}13^\circ$, porque a subida da linha média junto ao bordo
de fuga anula quase todo o efeito da curvatura.

Na zona linear, o método dos painéis prevê o maior declive (6,78~rad$^{-1}$)
por incluir a espessura. O CFD (6,26~rad$^{-1}$) fica praticamente igual à
teoria dos perfis delgados ($2\pi$), porque o aumento devido à espessura é
compensado pela redução devida à viscosidade.

O CFD é linear até cerca de $10^\circ$ e o perfil entra em perda entre $15^\circ$
e $20^\circ$. A $20^\circ$ o escoamento está descolado em quase todo o
extradorso e não é estacionário, o que os métodos potenciais, que admitem
escoamento sempre colado, não conseguem prever.

A $0^\circ$ e a $5^\circ$ as linhas de corrente e as distribuições de $C_p$ dos
painéis e do CFD são muito parecidas. A $20^\circ$ os painéis preveem um pico de
sucção de $-21{,}8$, enquanto o CFD mostra um pico muito menor e um patamar de
pressão no extradorso, típico do escoamento descolado.

Os métodos potenciais calculam toda a curva $c_l(\alpha)$ em pouco mais de um
segundo e chegam para estimar $c_l$ e $\alpha_{L=0}$ na zona linear. O CFD
demora minutos por ângulo e exige muito mais preparação, mas é o único que
serve perto e depois da perda e para obter o arrasto.

\paragraph{Uso de inteligência artificial.} Utilizou-se IA (Claude, da
Anthropic) como apoio na programação em MATLAB (painelização, geração de figuras
e correção de código), na formatação do relatório em \LaTeX{} e na redação de
rascunhos do texto, que foram revistos pelo grupo.}
\tbd{confirmar com o grupo onde foi usada}
